% !TEX program = xelatex
\documentclass[UTF8,a4paper,zihao=-4,fontset=windows]{ctexart}
\usepackage[margin=3cm]{geometry}
\usepackage{booktabs,tabularx,longtable,enumitem,fancyhdr,needspace}
\usepackage{xurl}
\usepackage[hidelinks,unicode]{hyperref}
\setlength{\parindent}{2em}
\setlength{\parskip}{.3em}
\setlength{\emergencystretch}{3em}
\linespread{1.12}
\setlist{nosep,leftmargin=2em}
\pagestyle{fancy}\fancyhf{}\fancyfoot[C]{\thepage}\renewcommand{\headrulewidth}{0pt}
\hypersetup{pdftitle={AI工具使用详情},pdfauthor={},pdfsubject={四问论文初版的AI使用阶段记录}}
\begin{document}
\begin{center}{\LARGE\bfseries AI工具使用详情}\par\medskip
四问论文初版\quad 记录日期：2026年9月12日
\end{center}

\section{记录范围与当前状态}
本文件随《考虑变物性与径向收缩的圆柱药材干燥模型》初版提供，依据2026年人工智能工具使用规定整理。它记录已有材料能够证明的使用情况，不将程序检查、代理审阅或AI生成的意见写作参赛队员已经完成的人工审核。本版尚未取得覆盖全文的队员人工修改和核验记录；历史外部交付的精确模型型号、全部原始对话及人工决策过程亦未完全提供，以下明确列出这些缺口，不补造记录。

本次工作以已形成的四问代码、正式结果及独立复审为输入，整理成论文初稿；本轮没有从头重跑全部四问的历史数值求解，也没有把历史含潜热模型或尚未闭合的候选模型替换成现行答案。对是否采纳本稿及其模型解释，仍须由参赛队员作出独立判断，并记录实质修改。

\section{工具名称、版本与分工}
\begin{longtable}{p{.23\linewidth}p{.29\linewidth}p{.39\linewidth}}
\toprule 工具或记录标识&版本信息&实际用途及信息边界\\\midrule\endhead
Codex桌面会话及协作代理&当前会话角色说明为基于GPT-6；精确服务端构建号未提供&读取本地材料；审查公式和结论；生成LaTeX文字、图表脚本、附录与核验记录。协作代理的审阅属于AI辅助复核。\\
既有外部Pro交付&材料仅标识为Pro交付；精确模型名称和版本未完整记录&此前参与四问建模、程序与答案生成、定向复审；本版采用其中经项目核查后保留的结果。不能将订阅名称Pro等同于确定的模型版本。\\
程序工具链（非生成式AI）&Python、NumPy、SciPy、Matplotlib、TeX Live 2026与XeLaTeX；细版本见复现说明&执行数值计算、数据读取、作图、排版及PDF核查。程序成功运行不能单独证明模型物理正确。\\
\bottomrule
\end{longtable}

\section{可追溯的关键输入与交互}
下列为本次已取得对话中的原始要求节录或有明确标记的摘要，省去姓名、账号、私有绝对路径和无关内容；不是对历史全部对话的声称。

\subsection{整体复审任务：要求节录}
\begin{quote}
“请对……项目现在的四问结果做一次整体、批判性的模型复审。第四问完整解答和前三问收尾已经加入仓库，本机核查材料也已提供。”

“不要只检查文件齐不齐或重复认可原答案，也不要为了复杂而换模型。”

“现行主线为题给物性、等效传热传质、无显式潜热；二维用于检验中截面与全域终点，第四问加入附件2的收缩。该选择是复审起点，不是免于质疑的正确性结论。”

“不要直接用历史F或未闭合R\_L代替现行结果，不要用‘2--3天’倒拟合参数。”
\end{quote}
输入还明确：问题二读取新的完整逐秒结果；问题三、四执行时长为57.4741小时、51.0921小时；条件执行值、二维校正估计和严格数学界分开。任务要求检查边界口径、干质量、潜热、后段环境、降维范围、严格阈值与收缩归因。

\subsection{论文任务：原始要求}
\begin{quote}
“好的，你现在根据该项目中的关于本题的全部分析和结论，去完成初版论文的编写。要求：用latex，论文格式符合今年国赛官方要求。”
\end{quote}
AI据此读取现行官方格式和AI规定，采用电子稿摘要首页、无目录、匿名、连续页码、完整代码附录，并使用独立版本目录保留历史成果。

\subsection{协作过程摘要}
主代理负责统一假设、共同方程、适用范围、综合结论、引用协调、编译与交付核验。协作代理分别整理问题一和二、问题三和四、冻结数据对应的六张题表与五幅图，并对全文进行交叉阅读。代理明确收到的约束包括：只能引用取得的文件；不得伪称补算完成；不使用旧结果替代现行版本；不改写历史证据；区分论文结论和待补计算。

本次生成的图表以保存数组或已有图源为依据，不利用图像反读数字。第一、二问题表、第三、四问终点及域外空白由机器核对；论文编译后再次检查最终PDF的页数、题表、匿名信息和版面。具体实际通过项目以本版验收记录为准，不能把未执行的完整冷启动列为已完成。

\Needspace{10\baselineskip}
\section{输出采纳、修订与核验}
\begin{longtable}{p{.2\linewidth}p{.35\linewidth}p{.36\linewidth}}
\toprule 内容&初版的采纳方式&核验与限制\\\midrule\endhead
模型框架&采用题给物性、有效Robin边界及无显式潜热的条件模型；第四问为给定半径下的材料坐标模型&保留气固映射、真实密度与能量不闭合的讨论；不写作模型已获实验验证。\\
正式数值&第一问原现行结果；第二问新完整结果；第三、四问现行一维执行时长&沿用已有独立数值证据，本轮主要核对引述和图表。未重新冷启动全部历史求解。\\
二维结论&按问题、网格、时刻分别说明直接场与配对校正&问题四精细二维正式终点未完成直接补证；不将其写成连续严格保证。\\
环境情景&采用已有真实后段续算及选定加密结果&人为扰动不是实测误差或统计置信区间；第四问保持规定半径。\\
收缩效果&采用同附录4固定半径对照&约60.65\%属于给定几何和干基边界标度的共同条件效果，不叫纯几何效应或实际节能率。\\
论文与图表&AI起草LaTeX、表格、图注、附录和说明，再由代理交叉检查&队员尚未提供全文人工审核与修改记录；本项仍待实际完成并补记。\\
\bottomrule
\end{longtable}

已在AI协作阶段修订的具体例子包括：使问题四表题与实际列一致；统一材料导数记号；把气相浓度差的线性化与传质系数转换写清；用条件情景措辞替代不受证据支持的物理保证。此处的“修订”指AI和程序流程中的更改，不声称队员已经亲自验证。

\section{正式采用前需要补齐的真实记录}
\begin{enumerate}
\item 从实际使用界面或历史导出核实外部AI的工具、模型名称和版本；无法追溯的条目如实标注，不能按习惯补填。
\item 补充此前各问核心建模和求解中有代表性的真实输入、输出、采纳理由，以及队员主导核心工作的实际过程。不能以本文件中的AI摘要替代不存在的对话或人工证据。
\item 逐项人工核验题意解释、附录物性、四位表格、终点判据、参考文献、模型适用范围；记录实际修改内容、核验方法和结论。没有完成的事项继续标为未完成。
\item 根据实际采纳范围更新论文AI声明及本详情，检查论文和支撑材料匿名、完整与压缩大小要求。本版是可编辑初稿，不将上述未完成记录描述为已验收。
\end{enumerate}

\section{规范来源}
全国大学生数学建模竞赛组委会，《全国大学生数学建模竞赛论文格式规范（2026年修订稿）》，2026年3月3日：\url{https://www.mcm.edu.cn/html_cn/node/4cd596519c9eb9fbd866398f6df0caa3.html}。

全国大学生数学建模竞赛组委会，《全国大学生数学建模竞赛关于人工智能工具使用的规定（2026年试行）》，2026年8月3日，2026年9月1日起执行：\url{https://www.mcm.edu.cn/html_cn/node/fef94648f2836ab6cc81586f4c38512b.html}。

本版已读取上述官方网页及其附件；具体抓取记录保留在内部验收证据中。官方格式规范没有指定某个民间LaTeX模板，也没有强制统一字体、字号和行距。
\end{document}
